% ============================================================
% Hydration Leakage and Over-fetching in Next.js
% CAIS-2026 — Springer ACSAR (LNCS llncs class)
% ============================================================
\documentclass[runningheads]{llncs}

\usepackage[T1]{fontenc}
\usepackage[utf8]{inputenc}
\usepackage{graphicx}
\usepackage{listings}
\usepackage{xcolor}
\usepackage{hyperref}
\usepackage{booktabs}
\usepackage{array}
\usepackage{amsmath}
\usepackage{microtype}
\usepackage{url}
\usepackage{float}

\renewcommand\UrlFont{\color{blue}\rmfamily}

% ── Code listing style ────────────────────────────────────────
\definecolor{codebg}{HTML}{F8F8F8}
\definecolor{codestring}{HTML}{008000}
\definecolor{codecomment}{HTML}{696969}
\definecolor{codekeyword}{HTML}{0000FF}
\definecolor{alert}{HTML}{CC0000}

\lstdefinestyle{nextjs}{
  backgroundcolor=\color{codebg},
  basicstyle=\ttfamily\footnotesize,
  breakatwhitespace=false,
  breaklines=true,
  captionpos=b,
  commentstyle=\color{codecomment}\itshape,
  keywordstyle=\color{codekeyword}\bfseries,
  stringstyle=\color{codestring},
  numberstyle=\tiny\color{gray},
  numbers=left,
  numbersep=6pt,
  showspaces=false,
  showstringspaces=false,
  showtabs=false,
  tabsize=2,
  frame=single,
  rulecolor=\color{gray!40},
  xleftmargin=1.5em,
}

\lstdefinestyle{python}{
  style=nextjs,
  language=Python,
  morekeywords={dataclass,field,Any,None,True,False},
}

\lstdefinestyle{typescript}{
  style=nextjs,
  language=Java,
  morekeywords={export,interface,async,await,const,let,type,
                string,number,boolean,null,undefined,
                GetServerSideProps,Promise},
}

\lstdefinestyle{json}{
  style=nextjs,
  language={},
  morestring=[b]",
  stringstyle=\color{codestring},
}

% ── Document metadata ─────────────────────────────────────────
\begin{document}

\title{Hydration Leakage and Over-fetching in Next.js:\\
The Security Implications of Server-to-Client\\
Data Serialization}

\titlerunning{Hydration Leakage in Next.js}

\author{Arnav Garg}

\authorrunning{A. Garg}

\institute{
  Department of Computer Science\\
  \email{arnav@example.com}
}

\maketitle

% ─────────────────────────────────────────────────────────────
\begin{abstract}
Modern full-stack web frameworks such as Next.js achieve high performance
by pre-rendering pages on the server and then \emph{hydrating} them on
the client.  This hydration process requires serializing server-side data
into the HTML response so that the browser can reconstruct React's
component tree without an additional network round-trip.  We demonstrate
that this mechanism, when combined with the common developer anti-pattern
of returning raw database records as component props, causes sensitive
fields---including password hashes, API keys, Social Security Numbers, and
two-factor authentication secrets---to be silently embedded in plain-text
HTML and transmitted to every unauthenticated client.  The vulnerability
affects both the legacy Pages Router (via the \texttt{\_\_NEXT\_DATA\_\_}
script tag) and the modern App Router (via React Server Component streaming
payloads).  We present an intentionally vulnerable proof-of-concept
application, an automated Python scanner capable of detecting exposure
across both serialization channels, and a Data Transfer Object (DTO)
mitigation that eliminates 100\% of confirmed critical findings with
minimal code changes.  Our scanner confirmed 24 critical credential
exposures across two route implementations; after applying the DTO
pattern, zero critical findings remained.  We argue that explicit,
allowlist-based data stripping at the server--client boundary must become
a standard practice in Next.js development.
\end{abstract}

\keywords{Next.js \and React \and hydration \and data leakage \and
          serialization \and security \and DTO \and server-side rendering}

% ─────────────────────────────────────────────────────────────
\section{Introduction}
\label{sec:intro}

Server-Side Rendering (SSR) has undergone a significant renaissance since
the widespread adoption of React-based meta-frameworks.  Unlike traditional
single-page applications that fetch data entirely on the client,
SSR frameworks such as Next.js, Nuxt, and SvelteKit generate complete HTML
on the server and stream it to the browser, dramatically improving both
perceived performance and search-engine discoverability~\cite{next16docs}.

The process of making a server-rendered HTML page interactive---known as
\emph{hydration}---requires the client-side JavaScript runtime to receive
the same data that the server used to render the page.  In practice, this
data is embedded directly in the HTML response as a serialized JSON payload.
This design is intentional and necessary for React's reconciliation
algorithm to avoid redundant fetches, but it creates a silent security
hazard: \emph{every field in every object passed from the server to a
client component is publicly visible in the HTML source of the page.}

The problem is compounded by the intuitive but incorrect mental model many
developers hold.  Seeing that a component only \emph{displays} a user's
name and email, a developer may not realize that the entire user object---
including its \texttt{password\_hash}, \texttt{api\_key}, and
\texttt{ssn} fields---was serialized into the response.  No developer
tool warns about this.  No browser security policy blocks it.  The data
simply appears in plaintext in the page source, readable by any person or
automated script that requests the URL.

This paper makes the following contributions:

\begin{itemize}
  \item A systematic analysis of the two distinct serialization channels
        in Next.js through which sensitive data leaks
        (Section~\ref{sec:background}).
  \item An intentionally vulnerable proof-of-concept application
        demonstrating the attack on both the Pages Router and the App
        Router (Section~\ref{sec:methodology}).
  \item An open-source automated scanner that detects hydration leakage
        across both channels using field-name and value-pattern heuristics
        (Section~\ref{sec:methodology}).
  \item A Data Transfer Object mitigation strategy, its implementation,
        and empirical validation via scanner-confirmed zero findings
        (Section~\ref{sec:mitigation}).
  \item A discussion of scanner limitations, false positives, and
        directions for future tooling (Section~\ref{sec:discussion}).
\end{itemize}

% ─────────────────────────────────────────────────────────────
\section{Background and Related Work}
\label{sec:background}

\subsection{Server-Side Rendering and Hydration}

Next.js supports two rendering models relevant to this work.  In the
\emph{Pages Router} model, each page exports a \texttt{getServerSideProps}
function that runs on the server per request, fetches data, and returns it
as a plain JavaScript object.  Next.js then passes these props as
arguments to the React component and serializes the entire \texttt{props}
object into a JSON blob embedded in a \texttt{<script>} tag with
\texttt{id="\_\_NEXT\_DATA\_\_"}~\cite{next16docs}.  This tag is
transmitted to every client regardless of authentication state.

The \emph{App Router}, introduced in Next.js 13 and stabilised in version
14+, uses React Server Components (RSC)~\cite{rsc2023}.  Server
Components run exclusively on the server and can fetch data directly, but
any data passed as props to a \texttt{'use client'} Client Component must
cross the server--client serialization boundary.  The RSC protocol
encodes this data as a series of \texttt{self.\_\_next\_f.push([type,
payload])} inline script calls, where the payload contains a
newline-delimited stream of JSON-encoded React tree segments.

In both models, the mechanism is \emph{necessary} for correct hydration
but \emph{dangerous} when the data objects contain fields that were never
intended for client consumption.

\subsection{Relationship to Existing Vulnerability Classes}

Hydration leakage is distinct from classical web vulnerabilities in
several important ways.  Unlike Cross-Site Scripting
(XSS)~\cite{owasp_xss}, it does not require injecting malicious content;
the server itself is the origin of the exposure.  Unlike SQL
injection~\cite{owasp_sqli}, it does not involve a malformed query;
the application logic is otherwise correct.  Unlike information disclosure
via verbose error messages~\cite{cwe200}, the exposure is not a failure
mode but a side-effect of the normal, successful request--response cycle.

Content Security Policy (CSP) headers, which are the standard mitigation
for XSS, offer no protection: the leaking data is not injected script but
first-party server-rendered JSON~\cite{csp_mdn}.  HTTPS prevents
interception in transit, but the exposure is at the \emph{origin}---any
authenticated or unauthenticated browser request retrieves the data.
Rate limiting and bot detection can slow enumeration but cannot prevent
a single manual page view from exposing all serialized records.

\subsection{Prior Work}

Concerns about Next.js data exposure were first widely documented in
practitioner blog posts circa 2021--2022, following the release of
\texttt{getServerSideProps}~\cite{nwh2022}.  Netanel Basal~\cite{basal2022}
and the Next.js security documentation~\cite{next_security} warn that
\texttt{getServerSideProps} return values are client-visible, but neither
provides tooling for automated detection.

Academic treatment of framework-specific serialization risks remains
sparse.  Luo et al.~\cite{luo2023} analyze state exposure in SPAs through
Redux DevTools, which requires browser extension access.  Our work is
distinct in targeting the \emph{network-layer} exposure that requires no
special tooling---only an HTTP client---and in covering both the legacy and
modern Next.js routing paradigms.

The OWASP API Security Top 10~\cite{owasp_api2023} lists \emph{Excessive
Data Exposure} (API3:2023) as a top-tier risk, defining it as an API
returning more data than the client needs and relying on the client to
filter.  Our findings constitute a precise instantiation of this risk
class within the SSR hydration context.

% ─────────────────────────────────────────────────────────────
\section{Methodology}
\label{sec:methodology}

Our methodology proceeds in two phases: constructing an instrumentally
vulnerable application that faithfully models real-world developer
mistakes, and building an automated scanner to detect the resulting
exposure.

\subsection{Phase 1: Vulnerable Proof-of-Concept Application}

We implemented a Next.js 16.1.6 application using React 19 and
TypeScript, hosted on a local development server.  The application
contains a mock user database stored as a JSON file
(\texttt{src/db/users.json}) with three synthetic user records.  Each
record contains both publicly-appropriate fields and sensitive fields
that should never reach the client, as detailed in Table~\ref{tab:schema}.

\begin{table}[H]
  \caption{Mock user database schema showing public and sensitive fields.}
  \label{tab:schema}
  \centering
  \begin{tabular}{lll}
    \toprule
    \textbf{Field} & \textbf{Type} & \textbf{Classification} \\
    \midrule
    \texttt{id}               & integer  & Public \\
    \texttt{name}             & string   & Public \\
    \texttt{email}            & string   & Public \\
    \texttt{role}             & string   & Public \\
    \texttt{created\_at}      & ISO-8601 & Public \\
    \midrule
    \texttt{password\_hash}   & string   & \textbf{CRITICAL — Bcrypt hash} \\
    \texttt{api\_key}         & string   & \textbf{CRITICAL — Live API key} \\
    \texttt{ssn}              & string   & \textbf{CRITICAL — Social Security Number} \\
    \texttt{two\_factor\_secret} & string & \textbf{CRITICAL — TOTP seed} \\
    \bottomrule
  \end{tabular}
\end{table}

\subsubsection{Vulnerability 1: Pages Router (\texttt{\_\_NEXT\_DATA\_\_})}

The Pages Router page at \texttt{/vulnerable-pages} exports
\texttt{getServerSideProps}, which fetches all user records from the
database and returns them as the \texttt{users} prop.  The component
renders only the user's name, email, and role.  Listing~\ref{lst:vuln_pages}
shows the vulnerable pattern.

\begin{lstlisting}[style=typescript, caption={Vulnerable Pages Router: raw
  database objects returned as props.}, label={lst:vuln_pages}]
// src/pages/vulnerable-pages.tsx
export const getServerSideProps: GetServerSideProps<Props> = async () => {
  const users = await getUsers(); // returns full User[] from DB

  // VULNERABILITY: raw DB objects returned directly.
  // Next.js serializes ALL fields into __NEXT_DATA__ in the HTML.
  return { props: { users } };
};
\end{lstlisting}

The resulting HTML contains the following embedded JSON, which is
accessible to any HTTP client that requests the page:

\begin{lstlisting}[style=json, caption={Excerpt of \_\_NEXT\_DATA\_\_ in
  the HTML response, showing leaked credentials.}, label={lst:next_data}]
<script id="__NEXT_DATA__" type="application/json">
{
  "props": {
    "pageProps": {
      "users": [{
        "id": 1,
        "name": "Alice Admin",
        "email": "alice@example.com",
        "password_hash": "$2b$12$EixZaYVK1fsbw1ZfbX3OXe...",
        "two_factor_secret": "JBSWY3DPEHPK3PXP",
        "api_key": "sk_live_1234567890abcdef",
        "ssn": "000-00-0001",
        "role": "admin"
      }, ...]
    }
  }
}
</script>
\end{lstlisting}

\subsubsection{Vulnerability 2: App Router (RSC Payload)}

The App Router page at \texttt{/vulnerable} uses a Server Component to
fetch users and pass the raw array to a \texttt{'use client'} Client
Component.  Listing~\ref{lst:vuln_app} shows the Server Component.

\begin{lstlisting}[style=typescript, caption={Vulnerable App Router:
  raw objects crossing the RSC serialization boundary.},
  label={lst:vuln_app}]
// src/app/vulnerable/page.tsx
export default async function VulnerableAppRouterPage() {
  const users = await getUsers(); // full User[] from DB

  // VULNERABILITY: entire User[] passed to Client Component.
  // Next.js encodes all fields into self.__next_f.push() payloads.
  return <UserListClient users={users} />;
}
\end{lstlisting}

Although the \texttt{UserListClient} component only renders
\texttt{user.name} and \texttt{user.email} in its JSX, the entire
\texttt{User} object---including \texttt{password\_hash},
\texttt{api\_key}, \texttt{ssn}, and \texttt{two\_factor\_secret}---is
serialized into the RSC streaming payload embedded in the initial HTML
response.

The RSC protocol encodes data as consecutive inline script calls of the
form \texttt{self.\_\_next\_f.push([type, data])}, where row type 1
carries JSON-encoded component tree segments.  Unlike the Pages Router
\texttt{\_\_NEXT\_DATA\_\_} blob, the RSC payload is fragmented across
many such calls and requires parsing of a newline-delimited format where
each line is \texttt{<row\_id>:<json\_value>}.

\subsection{Phase 2: Automated Vulnerability Scanner}

We implemented a Python 3.12 scanner (\texttt{scanner.py}, $\approx$310
lines) that automates the detection of hydration leakage.  The scanner
uses \texttt{requests} for HTTP fetching and \texttt{BeautifulSoup} for
HTML parsing.

\subsubsection{Extraction Pipeline}

The scanner implements two extraction paths corresponding to the two
Next.js serialization channels:

\begin{enumerate}
  \item \textbf{Pages Router}: BeautifulSoup locates the
        \texttt{<script id="\_\_NEXT\_DATA\_\_">} tag and parses its
        \texttt{string} attribute as JSON.

  \item \textbf{App Router}: A compiled regular expression
        \texttt{r'self\textbackslash .\_\_next\_f\textbackslash
        .push\textbackslash(\textbackslash s*(\textbackslash[
        \textbackslash s\textbackslash S]*?\textbackslash])\textbackslash
        s*\textbackslash)'} finds all RSC push calls in the HTML.  Each
        matched payload is parsed, the row type inspected, and row-type-1
        payloads further split by newline and parsed as individual JSON
        objects.  All recovered objects are collected into a list for
        further analysis.
\end{enumerate}

\subsubsection{Recursive Tree Walker}

Once JSON data is recovered from either channel, the scanner recursively
walks the Python object graph.  At each dictionary key, two independent
checks run:

\begin{enumerate}
  \item \textbf{Field-name pattern matching}: The key is tested against
        29 compiled regular expressions organised into three severity
        tiers.  The patterns are ordered from most specific to most
        general (e.g., \texttt{password\_hash} before
        \texttt{password}) to ensure the most descriptive label is
        assigned when a key matches multiple patterns.  The first matching
        pattern wins; a \texttt{set} of already-flagged paths is
        maintained across the recursion to prevent duplicate findings.

  \item \textbf{Value pattern matching}: If the field was not flagged by
        a name pattern and its value is a string, the value is tested
        against eight compiled value-format regular expressions covering
        bcrypt hashes, Stripe-format API keys, SSN format, JWTs, GitHub
        personal access tokens, AWS access key IDs, Base32 strings
        (potential TOTP secrets), and 13--16-digit numeric sequences
        (potential credit card numbers).
\end{enumerate}

Table~\ref{tab:patterns} summarises the field-name pattern coverage by
severity tier.

\begin{table}[H]
  \caption{Scanner field-name patterns by severity tier.}
  \label{tab:patterns}
  \centering
  \begin{tabular}{p{1.8cm}p{6.5cm}p{2cm}}
    \toprule
    \textbf{Severity} & \textbf{Field name patterns (regex)} & \textbf{Count} \\
    \midrule
    CRITICAL &
      \texttt{password\_hash}, \texttt{hashed\_password},
      \texttt{password}, \texttt{passwd},
      \texttt{two\_factor\_secret}, \texttt{totp\_secret},
      \texttt{otp\_secret}, \texttt{private\_key},
      \texttt{api\_key}, \texttt{apikey},
      \texttt{social\_security}, \texttt{ssn},
      \texttt{credit\_card}, \texttt{card\_number}, \texttt{cvv},
      \texttt{secret}
    & 16 \\
    \midrule
    HIGH &
      \texttt{access\_token}, \texttt{refresh\_token},
      \texttt{auth\_token}, \texttt{session\_token},
      \texttt{\textbackslash btoken\textbackslash b}, \texttt{jwt},
      \texttt{bearer}, \texttt{encryption\_key}, \texttt{signing\_key}
    & 9 \\
    \midrule
    MEDIUM &
      \texttt{webhook\_secret}, \texttt{internal\_id},
      \texttt{is\_admin}, \texttt{admin}, \texttt{stripe\_customer}
    & 5 \\
    \bottomrule
  \end{tabular}
\end{table}

\subsubsection{Findings Store}

Every scan run is automatically persisted to a \texttt{findings/}
directory as a timestamped JSON file.  A \texttt{findings/index.json}
catalog is maintained and updated after each run, recording the
timestamp, user-supplied label, target count, and per-severity finding
counts.  This enables longitudinal comparison across scan runs, as
exploited in Section~\ref{sec:results}.

% ─────────────────────────────────────────────────────────────
\section{Results}
\label{sec:results}

\subsection{Baseline Scan: Vulnerable Routes}

Running the scanner against both vulnerable routes with the command:

\begin{lstlisting}[style=python, numbers=none]
python3 scanner.py --targets targets.txt --label vulnerable-baseline
\end{lstlisting}

produced the results in Table~\ref{tab:results_vuln}.

\begin{table}[H]
  \caption{Scanner findings on vulnerable routes (pre-mitigation).}
  \label{tab:results_vuln}
  \centering
  \begin{tabular}{llrrr}
    \toprule
    \textbf{Route} & \textbf{Channel} & \textbf{CRITICAL} & \textbf{HIGH} & \textbf{MEDIUM} \\
    \midrule
    \texttt{/vulnerable-pages} & \texttt{\_\_NEXT\_DATA\_\_} & 12 & 4$^{*}$ & 0 \\
    \texttt{/vulnerable}       & RSC payload                  & 12 & 40$^{*}$ & 0 \\
    \midrule
    \textbf{Total}             &                              & \textbf{24} & 44 & 0 \\
    \bottomrule
  \end{tabular}

  \smallskip
  \footnotesize{$^{*}$ All HIGH findings are false positives (see
    Section~\ref{sec:discussion}).}
\end{table}

The 24 confirmed CRITICAL findings decompose exactly as expected: four
sensitive fields (\texttt{password\_hash}, \texttt{api\_key},
\texttt{ssn}, \texttt{two\_factor\_secret}) per user, across three users,
across two routes ($4 \times 3 \times 2 = 24$).  Every finding was
confirmed by manual inspection of the raw HTML.  Listing~\ref{lst:finding}
shows a representative scanner output entry.

\begin{lstlisting}[style=json, caption={Representative CRITICAL finding
  from the scanner JSON report.}, label={lst:finding}]
{
  "field_path": "__NEXT_DATA__.props.pageProps.users[0].password_hash",
  "field_name": "password_hash",
  "value_preview": "$2b$12...RQYq  [60 chars]",
  "reason": "Password hash -- field name matches /password_hash/i",
  "severity": "CRITICAL",
  "channel": "__NEXT_DATA__",
  "match_type": "field_name"
}
\end{lstlisting}

The dot-notation \texttt{field\_path} precisely locates each leaked field
within the serialized JSON tree, providing actionable remediation
information to a developer.

\subsection{Post-Mitigation Scan: Secure Routes}

After implementing the DTO mitigation described in
Section~\ref{sec:mitigation}, the scanner was run against both the
vulnerable and the mitigated routes simultaneously
(Table~\ref{tab:results_comparison}).

\begin{table}[H]
  \caption{Scanner findings before and after DTO mitigation.}
  \label{tab:results_comparison}
  \centering
  \begin{tabular}{llrrr}
    \toprule
    \textbf{Route} & \textbf{State} & \textbf{CRITICAL} & \textbf{HIGH$^{*}$} & \textbf{MEDIUM} \\
    \midrule
    \texttt{/vulnerable-pages} & Baseline   & 12 & 4  & 0 \\
    \texttt{/vulnerable}       & Baseline   & 12 & 40 & 0 \\
    \texttt{/secure-pages}     & Mitigated  & \textbf{0}  & 4  & 0 \\
    \texttt{/secure}           & Mitigated  & \textbf{0}  & 36 & 0 \\
    \bottomrule
  \end{tabular}

  \smallskip
  \footnotesize{$^{*}$ All HIGH findings on mitigated routes are false
    positives (see Section~\ref{sec:discussion}). The reduction from 40
    to 36 HIGH on the App Router reflects four fewer \texttt{users[n].role}
    false-positive entries, as the RSC tree structure differs slightly
    between routes.}
\end{table}

The DTO mitigation eliminated 100\% of confirmed CRITICAL findings.  The
zero-finding result on the mitigated routes was further validated by
direct inspection of the HTML source: no \texttt{password\_hash},
\texttt{api\_key}, \texttt{ssn}, or \texttt{two\_factor\_secret} field
appeared anywhere in the response payload.

% ─────────────────────────────────────────────────────────────
\section{Mitigation: The Data Transfer Object Pattern}
\label{sec:mitigation}

\subsection{Design Rationale}

The root cause of hydration leakage is the absence of an explicit
boundary between the server's internal data representation and the data
contract exposed to clients.  Two approaches to enforcing such a boundary
are possible:

\begin{enumerate}
  \item \textbf{Blocklist (field deletion)}: After fetching, explicitly
        delete sensitive fields before returning props, e.g.,
        \texttt{delete user.password\_hash}.

  \item \textbf{Allowlist (DTO projection)}: Define a separate type
        containing only the fields that are permitted to cross the
        boundary, and construct instances of that type from the database
        records.
\end{enumerate}

We adopt and recommend the \emph{allowlist} approach.  A blocklist is
\emph{open by default}---when a new sensitive column is added to the
database schema, it leaks to clients until a developer remembers to add
it to the deletion list.  An allowlist is \emph{closed by default}---new
fields are invisible to clients until explicitly declared in the DTO, and
the TypeScript type system enforces this at compile time.

\subsection{Implementation}

We defined a \texttt{PublicUser} DTO interface in
\texttt{src/lib/dto.ts} containing only the five fields appropriate for
client consumption.  Listing~\ref{lst:dto} shows the full implementation.

\begin{lstlisting}[style=typescript, caption={DTO definition and
  transformation function in \texttt{src/lib/dto.ts}.},
  label={lst:dto}]
export interface PublicUser {
  id:         number;
  name:       string;
  email:      string;
  role:       string;
  created_at: string;
  // password_hash, api_key, ssn, two_factor_secret: absent.
}

export function toPublicUser(user: User): PublicUser {
  return {
    id:         user.id,
    name:       user.name,
    email:      user.email,
    role:       user.role,
    created_at: user.created_at,
  };
}

export const toPublicUsers = (users: User[]): PublicUser[] =>
  users.map(toPublicUser);
\end{lstlisting}

The explicit destructuring in \texttt{toPublicUser} ensures that
TypeScript enforces the shape at compile time: any field not listed in
the return literal is structurally excluded.  Furthermore, because Client
Components in the secure route accept \texttt{PublicUser[]} rather than
\texttt{User[]}, any accidental reference to \texttt{user.password\_hash}
within client code is a compile-time type error, not a runtime leak.

\subsubsection{Pages Router Integration}

The secure Pages Router page \texttt{/secure-pages} applies
\texttt{toPublicUsers} before returning props:

\begin{lstlisting}[style=typescript, caption={DTO-secured
  \texttt{getServerSideProps}.}, label={lst:secure_pages}]
export const getServerSideProps: GetServerSideProps<Props> =
  async () => {
    const users = await getUsers();
    return { props: { users: toPublicUsers(users) } };
  };
\end{lstlisting}

Only the \texttt{PublicUser} shape now enters the
\texttt{\_\_NEXT\_DATA\_\_} serialization pipeline.

\subsubsection{App Router Integration}

The secure App Router Server Component \texttt{/secure} applies the
transformation before passing data to the Client Component:

\begin{lstlisting}[style=typescript, caption={DTO-secured
  App Router Server Component.}, label={lst:secure_app}]
export default async function SecureAppRouterPage() {
  const users      = await getUsers();
  const publicUsers = toPublicUsers(users); // strip before boundary

  return <SecureUserListClient users={publicUsers} />;
}
\end{lstlisting}

The RSC serialization boundary only ever observes
\texttt{PublicUser} objects.  Sensitive fields are removed before they
can enter the \texttt{self.\_\_next\_f.push} payload.

\subsection{Code Change Footprint}

Table~\ref{tab:changefootprint} quantifies the implementation cost of the
mitigation.  The DTO module itself requires approximately 20 lines of
TypeScript.  Each route requires a single one-line change to wrap the
database result in \texttt{toPublicUsers()}.

\begin{table}[H]
  \caption{Lines of code changed to apply the DTO mitigation per route.}
  \label{tab:changefootprint}
  \centering
  \begin{tabular}{lrl}
    \toprule
    \textbf{File} & \textbf{Lines} & \textbf{Change} \\
    \midrule
    \texttt{src/lib/dto.ts}               & 20 & New file: DTO interface and \texttt{toPublicUser} \\
    \texttt{src/pages/secure-pages.tsx}   & 1  & Wrap \texttt{getUsers()} result with \texttt{toPublicUsers()} \\
    \texttt{src/app/secure/page.tsx}      & 1  & Wrap \texttt{getUsers()} result with \texttt{toPublicUsers()} \\
    \midrule
    \textbf{Total}                        & 22 & Eliminates 24 CRITICAL findings \\
    \bottomrule
  \end{tabular}
\end{table}

% ─────────────────────────────────────────────────────────────
\section{Discussion and Limitations}
\label{sec:discussion}

\subsection{False Positives in the Scanner}

The scanner produced HIGH-severity false positives on all scanned routes,
arising from two structural properties of the Next.js RSC payload and
Pages Router data.

\subsubsection{Base32 Value Pattern}

The scanner's value heuristic for TOTP secrets uses the pattern
\texttt{\textasciicircum[A-Z2-7]+=*\$}, which matches any string composed
entirely of Base32 alphabet characters.  Short English words and React
internal values frequently satisfy this constraint: \texttt{``user''},
\texttt{``admin''}, \texttt{``flex''}, \texttt{``center''}, \texttt{``en''}
(the HTML \texttt{lang} attribute), and
\texttt{``development''} (the Next.js \texttt{buildId}) all match.  These
are React component tree metadata, HTML attributes, and CSS class names,
not TOTP secrets.

\subsubsection{Field Name Substring Matching: \texttt{className}}

React's RSC payload encodes component props as JSON, including the
\texttt{className} prop used for CSS class assignment.  The field name
\texttt{className} contains the substring \texttt{ssN}
(case-insensitively), triggering the scanner's SSN field-name pattern.
The associated values are Tailwind CSS utility strings such as
\texttt{``p-8 font-sans''} and Geist font identifier strings, not Social
Security Numbers.

\subsubsection{Mitigating False Positives}

Both false positive classes are structurally distinguishable from genuine
findings: they occur at paths containing well-known React-internal field
names (\texttt{className}, \texttt{lang}, \texttt{rel},
\texttt{parallelRouterKey}, \texttt{style.*}) and are present in both
vulnerable and mitigated routes equally.  A production-grade scanner
could eliminate them by maintaining an allowlist of internal RSC field
names to exclude from pattern matching, or by applying a minimum-length
filter on value pattern matches.  We leave this refinement to future work.

Importantly, the false positive rate on CRITICAL findings from genuine
user-data paths is \textbf{zero}: every CRITICAL finding that references
a user array path (e.g.,
\texttt{users[0].password\_hash}) was confirmed to be a true positive.

\subsection{Scope and Generalizability}

Our proof-of-concept uses a local JSON file as the mock database.
Production Next.js applications typically use relational or document
databases via ORMs such as Prisma~\cite{prisma}, which return typed model
objects with even more fields than our minimal schema.  The hydration
leakage mechanism is identical; ORMs may in fact \emph{increase} risk by
returning comprehensive record objects by default.

The vulnerability is not specific to Next.js.  Any SSR framework that
serializes server-fetched data into the HTML response and passes it to
client-side components---including Nuxt 3 (Vue), SvelteKit, Remix, and
Astro in hybrid mode---is susceptible to the same class of exposure when
developers pass raw database records across the server--client boundary.
The DTO pattern we describe is framework-agnostic.

\subsection{Developer Tooling Gap}

No current build-time or lint-time tooling warns developers when a
\texttt{getServerSideProps} return value or a Server Component prop
contains fields that were not declared as intentionally public.
TypeScript's structural type system cannot distinguish ``this interface
is a database record'' from ``this interface is a safe DTO''; without
nominal typing or explicit boundary annotations, the compiler permits
the leak silently.

An ESLint plugin that enforces DTO use at \texttt{getServerSideProps}
return sites and at RSC-to-Client-Component prop boundaries would close
this tooling gap.  We propose this as a direction for future work.

% ─────────────────────────────────────────────────────────────
\section{Conclusion}
\label{sec:conclusion}

We have demonstrated that the hydration mechanism in Next.js---common to
both the Pages Router and the App Router---constitutes a latent
information-disclosure vulnerability when developers pass raw database
records as component props.  The exposure requires no attacker privilege,
no network interception, and no browser extension: the data is present in
the plain-text HTML of every response.

Our automated scanner confirmed 24 CRITICAL credential exposures across
two route implementations in our proof-of-concept application.  These
exposures included bcrypt password hashes, live API keys, Social Security
Numbers, and TOTP two-factor authentication seeds.  After applying the
Data Transfer Object mitigation---a structural pattern that allowlists
permitted fields before they enter the serialization pipeline---the scanner
confirmed zero CRITICAL findings.  The complete mitigation required
approximately 22 lines of TypeScript.

We recommend three practices for Next.js developers and security
reviewers:

\begin{enumerate}
  \item \textbf{Adopt the DTO allowlist pattern universally.} Every data
        access in \texttt{getServerSideProps} and every prop passed from
        a Server Component to a Client Component should be projected
        through an explicit DTO before crossing the boundary.

  \item \textbf{Treat the serialized HTML payload as a public API
        surface.} Any field that would be inappropriate in a REST
        response to an unauthenticated caller is also inappropriate in
        the hydration payload.

  \item \textbf{Include hydration payload review in security audits.}
        Existing scanning tools focus on XSS, injection, and
        authentication bypasses.  Automated checks against
        \texttt{\_\_NEXT\_DATA\_\_} and RSC push payloads should be added
        to standard web-application security assessment methodology.
\end{enumerate}

As React-based SSR frameworks continue to proliferate and become the
default choice for full-stack web development, the boundary between
server-side data and client-side rendering deserves the same rigorous
security scrutiny historically applied to API design.

% ─────────────────────────────────────────────────────────────
\bibliographystyle{splncs04}
\bibliography{references}

\end{document}
